% !TITLE 赠卫八处士
% !AUTHOR 杜甫
% !DYNASTY 唐
% !ORDER 24

\begin{poem}
人生不相见，动如参与商\textsuperscript{1}。
今夕复何夕\textsuperscript{2}，共此灯烛光。
少壮能几时？鬓发各已苍\textsuperscript{3}！
访旧半为鬼\textsuperscript{4}，惊呼热中肠\textsuperscript{5}。
焉知二十载，重上君子堂。
昔别君未婚，儿女忽成行\textsuperscript{6}。
怡然敬父执\textsuperscript{7}，问我来何方。
问答乃未已，驱儿罗酒浆\textsuperscript{8}。
夜雨剪春韭，新炊间黄粱\textsuperscript{9}。
主称会面难，一举累十觞\textsuperscript{10}。
十觞亦不醉，感子故意长\textsuperscript{11}。
明日隔山岳，世事两茫茫。
\end{poem}

\begin{annotations}
\notetext{1}{参与商：参星（shēn）和商星，一颗在东，一颗在西，永不同时出现在天空。比喻人与人的离别——有些人一别就再也见不到了。}
\notetext{2}{今夕复何夕：今夜是什么夜晚啊（竟能与你相聚）。“复何夕”是感叹语——这样美好的夜晚人生能有几次？}
\notetext{3}{苍：灰白色。两人鬓发都已花白——从少年变成了老年。}
\notetext{4}{访旧半为鬼：打听旧友的下落，一半已经去世了。}
\notetext{5}{热中肠：心中火辣辣地痛。形容听到噩耗时的震撼感受。}
\notetext{6}{成行：排成行列。当年分别时还没结婚，现在孩子已经成排成行了——写时间流逝之快。}
\notetext{7}{父执：父亲的朋友。孩子们恭敬地称呼父亲的老朋友。}
\notetext{8}{罗酒浆：摆上酒和菜。罗，排列、摆出。}
\notetext{9}{间黄粱：饭里掺着黄粱（小米的一种）。夜雨中剪了春韭，新煮了黄粱饭——普通人家的家常饭菜。}
\notetext{10}{一举累十觞：一举杯就喝了十杯。累，累积；觞（shāng），酒杯。}
\notetext{11}{故意长：旧日的情意深长。不是酒不醉人，是感情比酒更浓。}
\end{annotations}

\begin{appreciation}
《赠卫八处士》是中国文学史上最温暖的重逢诗之一。759年，杜甫被贬官，赴任路上绕道去看二十年没见的老朋友卫八，老朋友留他吃了一顿饭。全诗没有一句大话，只是一顿饭，一盏灯，几杯酒。

人生不相见，动如参与商。参星和商星，一颗在东，一颗在西，此出彼没，永不同时挂在天上。拿天文规律比喻离别：有些人一别，就真是永别。所以下一句格外珍贵：今夕复何夕，共此灯烛光——今夜是什么日子，竟能和你坐在同一盏灯下。

少壮能几时？鬓发各已苍。苍是灰白。二十年，两个少年都白了头。访旧半为鬼，惊呼热中肠——问起老朋友的下落，一半已经不在人世。热中肠，心口发烫，是听到噩耗时最真实的生理反应。焉知二十载，重上君子堂——哪能想到，还能重新走进你的厅堂。

昔别君未婚，儿女忽成行。当年分别时你还没结婚，如今孩子都成排了。一个忽字——时间不是慢慢走的，是跳过来的。孩子们怡然敬父执，问我来何方——父执是父亲的朋友，他们恭恭敬敬地喊伯伯。

夜雨剪春韭，新炊间黄粱。全诗最好的两句：雨夜里剪下的韭菜，新煮的、掺了小米的饭。不是什么好菜，是一个穷朋友能把家底端上来的情分。主称会面难，一举累十觞。十觞亦不醉，感子故意长——情意比酒长。

明日隔山岳，世事两茫茫。烛光、春韭、十杯酒，到这里都显得薄。刚重逢，又要离别；刚暖起来，又要凉下去。杜甫没有哭喊，只是把事实说了一遍。

还记得《行行重行行》和《生年不满百》吗？一个写"相去万余里"的相思，一个说人生太短；到杜甫这里，二十年的离散熬成了一夜重逢——正因为人生太短，这一夜才这么重。

你将来也会有一个二十年不见的老朋友，再见面，大概也是一顿饭。好朋友要常约饭。你以后回来看我，不用带东西——我给你剪春韭。
\end{appreciation}
